% ══════════════════════════════════════════════════════════════════════════════
%                         فصل دهم
%                 لیبرالیسم کلاسیک و نئولیبرالیسم
% ══════════════════════════════════════════════════════════════════════════════

\chapter{لیبرالیسم کلاسیک و نئولیبرالیسم}
\label{ch:liberalism}

\begin{flushright}
\textit{«تنها آزادی‌ای که شایسته این نام است، پیگیری خیر خودمان به شیوه خودمان است.»}

— جان استوارت میل
\end{flushright}

% ══════════════════════════════════════════════════════════════════════════════
%                         مقدمه فصل
% ══════════════════════════════════════════════════════════════════════════════

\lettrine[lines=3, lraise=0.1, nindent=0.5em]{\textcolor{CenterGreen}{ل}}{یبرالیسم} 
شاید موفق‌ترین ایدئولوژی تاریخ مدرن باشد. حقوق بشر، دموکراسی، آزادی بیان، حکومت قانون، بازار آزاد—همه این‌ها از سنت لیبرالی برآمده‌اند. اما لیبرالیسم واژه‌ای گیج‌کننده است: در آمریکا، «لیبرال» یعنی چپ؛ در اروپا، اغلب یعنی راست اقتصادی. در این فصل، این سردرگمی را روشن می‌کنیم و سیر تحول لیبرالیسم را از لاک تا نئولیبرالیسم دنبال می‌کنیم.

% ══════════════════════════════════════════════════════════════════════════════
\section{لیبرالیسم چیست؟}
% ══════════════════════════════════════════════════════════════════════════════

\subsection{سردرگمی معنایی}

\begin{warningbox}[هشدار: لیبرال در آمریکا و اروپا]
\textbf{در آمریکا:}
\begin{itemize}[nosep]
    \item «لیبرال» = چپ میانه، دموکرات
    \item طرفدار دولت رفاه و حقوق اقلیت‌ها
    \item مخالف محافظه‌کاران
\end{itemize}

\textbf{در اروپا:}
\begin{itemize}[nosep]
    \item «لیبرال» = طرفدار بازار آزاد، دولت کوچک
    \item اغلب راست میانه اقتصادی
    \item ممکن است در مسائل اجتماعی لیبرال باشد
\end{itemize}

\textbf{علت:} در آمریکا، چپ میانه واژه «سوسیالیسم» را کنار گذاشت و «لیبرال» را برگزید. در اروپا، سوسیالیست‌ها و سوسیال‌دموکرات‌ها خود نام‌گذاری کردند.
\end{warningbox}

\subsection{ویژگی‌های مشترک لیبرالیسم}

\begin{defbox}[تعریف: لیبرالیسم (هسته مشترک)]
با وجود تنوع، اغلب لیبرال‌ها در این اصول اشتراک دارند:
\begin{itemize}[nosep]
    \item \textbf{فردگرایی:} فرد واحد اصلی اخلاق و سیاست است
    \item \textbf{آزادی:} حق فرد برای زندگی به شیوه خود
    \item \textbf{برابری:} برابری در برابر قانون و حقوق پایه
    \item \textbf{عقل‌گرایی:} باور به توانایی عقل در حل مسائل
    \item \textbf{پیشرفت:} جامعه می‌تواند بهتر شود
    \item \textbf{رضایت:} حکومت مشروع بر رضایت مردم استوار است
    \item \textbf{حکومت محدود:} قدرت دولت باید محدود باشد
\end{itemize}
\end{defbox}

\begin{table}[htbp]
\centering
\caption{انواع لیبرالیسم}
\label{tab:liberalism-types}
\begin{tabular}{r|c|c|c}
\toprule
\textbf{نوع} & \textbf{دوره} & \textbf{تأکید} & \textbf{نمونه} \\
\midrule
کلاسیک & قرن ۱۷-۱۹ & آزادی منفی، بازار & لاک، اسمیت \\
اجتماعی & اواخر ۱۹ & آزادی مثبت، رفاه & گرین، دیویی \\
نئولیبرالیسم & ۱۹۷۰s-امروز & بازگشت به بازار & هایک، فریدمن \\
برابری‌خواه & معاصر & عدالت توزیعی & رالز \\
\bottomrule
\end{tabular}
\end{table}

% ══════════════════════════════════════════════════════════════════════════════
\section{لیبرالیسم کلاسیک: ریشه‌ها}
% ══════════════════════════════════════════════════════════════════════════════

\subsection{جان لاک: پدر لیبرالیسم}

\begin{personbox}[جان لاک (۱۶۳۲-۱۷۰۴)]

\textbf{ملیت:} انگلیسی

\textbf{آثار کلیدی:} \book{دو رساله در باب حکومت} (۱۶۸۹)، \book{رساله در باب فهم بشری}، \book{نامه‌ای در باب تسامح}

\textbf{ایده محوری:} حقوق طبیعی و حکومت محدود

\mediumspace

\person{لاک} بنیان‌گذار لیبرالیسم سیاسی مدرن است:

\textbf{حقوق طبیعی:}
\begin{itemize}[nosep]
    \item هر انسان حقوقی دارد که پیش از دولت وجود دارند
    \item \textbf{زندگی، آزادی، مالکیت}
    \item دولت برای حفاظت از این حقوق تأسیس می‌شود
    \item اگر دولت حقوق را نقض کند، مردم حق شورش دارند
\end{itemize}

\textbf{رضایت:}
\begin{itemize}[nosep]
    \item حکومت مشروع بر «رضایت محکومان» استوار است
    \item قرارداد اجتماعی: مردم بخشی از آزادی را واگذار می‌کنند
    \item اما حقوق اساسی قابل واگذاری نیست
\end{itemize}

\mediumspace

\textbf{نقل‌قول:} \textit{«هدف قانون، نه لغو یا محدود کردن، بلکه حفظ و گسترش آزادی است.»}
\end{personbox}

\begin{figure}[htbp]
\centering
\begin{tikzpicture}[scale=0.85]

% عنوان
\node[font=\large\bfseries, text=DarkGray] at (0, 7) {\fa{نظریه سیاسی لاک}};

% وضع طبیعی
\fill[CenterGreenLight, opacity=0.6, rounded corners=8pt] (-6, 3.5) rectangle (-2, 6);
\draw[line width=2pt, color=CenterGreen, rounded corners=8pt] (-6, 3.5) rectangle (-2, 6);
\node[text=CenterGreenDark, font=\bfseries] at (-4, 5.5) {\fa{وضع طبیعی}};
\node[text=DarkGray, font=\scriptsize, text width=3.5cm, align=center] at (-4, 4.5) {
    \fa{آزادی و برابری}\\
    \fa{حقوق طبیعی}\\
    \fa{اما ناامن}
};

% قرارداد
\draw[thickarrow, color=GoldAccent, line width=3pt] (-1.8, 4.75) -- (0.3, 4.75);
\node[text=GoldAccent, font=\small\bfseries] at (-0.75, 5.3) {\fa{قرارداد}};

% جامعه مدنی
\fill[RightBlueLight, opacity=0.6, rounded corners=8pt] (0.5, 3.5) rectangle (4.5, 6);
\draw[line width=2pt, color=RightBlue, rounded corners=8pt] (0.5, 3.5) rectangle (4.5, 6);
\node[text=RightBlueDark, font=\bfseries] at (2.5, 5.5) {\fa{جامعه مدنی}};
\node[text=DarkGray, font=\scriptsize, text width=3.5cm, align=center] at (2.5, 4.5) {
    \fa{حکومت محدود}\\
    \fa{حفاظت از حقوق}\\
    \fa{قانون}
};

% شرط
\fill[LeftRedLight, opacity=0.5, rounded corners=5pt] (-2, 1) rectangle (4, 2.5);
\draw[line width=1.5pt, color=LeftRed, rounded corners=5pt] (-2, 1) rectangle (4, 2.5);
\node[text=LeftRedDark, font=\small\bfseries] at (1, 2) {\fa{اگر دولت حقوق را نقض کند:}};
\node[text=DarkGray, font=\small] at (1, 1.4) {\fa{حق شورش!}};

% فلش
\draw[thickarrow, color=LeftRed] (2.5, 3.3) -- (2.5, 2.6);

\end{tikzpicture}
\caption{ساختار نظریه سیاسی جان لاک}
\label{fig:locke-theory}
\end{figure}

\subsection{آدام اسمیت: لیبرالیسم اقتصادی}

\begin{personbox}[آدام اسمیت (۱۷۲۳-۱۷۹۰)]

\textbf{ملیت:} اسکاتلندی

\textbf{آثار کلیدی:} \book{ثروت ملل} (۱۷۷۶)، \book{نظریه احساسات اخلاقی} (۱۷۵۹)

\textbf{ایده محوری:} دست نامرئی بازار

\mediumspace

\person{اسمیت} پدر اقتصاد مدرن و لیبرالیسم اقتصادی است:

\begin{itemize}[nosep]
    \item \textbf{دست نامرئی:} دنبال کردن نفع شخصی، به نفع عمومی می‌انجامد
    \item \textbf{تقسیم کار:} تخصص، بهره‌وری را افزایش می‌دهد
    \item \textbf{تجارت آزاد:} همه از تجارت سود می‌برند
    \item \textbf{ضد مرکانتیلیسم:} دولت نباید در اقتصاد دخالت کند
\end{itemize}

\mediumspace

\textbf{اما:} اسمیت «بنیادگرای بازار» نبود:
\begin{itemize}[nosep]
    \item از آموزش عمومی دفاع کرد
    \item از مالیات تصاعدی حمایت کرد
    \item انحصارها را نقد کرد
    \item از زیرساخت‌های دولتی دفاع کرد
\end{itemize}
\end{personbox}

\begin{quotebox}[آدام اسمیت، ثروت ملل]
\textit{«این از خیرخواهی قصاب، آبجوساز یا نانوا نیست که ما انتظار شام خود را داریم، بلکه از توجه آنان به منفعت خودشان است.»}
\end{quotebox}

\subsection{جان استوارت میل: لیبرالیسم و آزادی}

\begin{personbox}[جان استوارت میل (۱۸۰۶-۱۸۷۳)]

\textbf{ملیت:} انگلیسی

\textbf{آثار کلیدی:} \book{درباره آزادی} (۱۸۵۹)، \book{فایده‌گرایی}، \book{انقیاد زنان}

\textbf{ایده محوری:} اصل آسیب و آزادی فردی

\mediumspace

\person{میل} لیبرالیسم را به اوج فلسفی رساند:

\textbf{اصل آسیب:}
\begin{itemize}[nosep]
    \item تنها توجیه مشروع برای محدود کردن آزادی فرد، \textbf{جلوگیری از آسیب به دیگران} است
    \item «خیر خود فرد» کافی نیست—پدرسالاری نادرست است
    \item جامعه حق ندارد عقاید یا سبک زندگی را تحمیل کند
\end{itemize}

\textbf{آزادی بیان:}
\begin{itemize}[nosep]
    \item حتی عقاید غلط باید آزاد باشند
    \item حقیقت از برخورد آرا بیرون می‌آید
    \item سانسور، حتی اگر نظر غلط را سانسور کند، زیان‌بار است
\end{itemize}

\textbf{فمینیسم:}
\begin{itemize}[nosep]
    \item از حق رأی زنان دفاع کرد
    \item انقیاد زنان را با اصول لیبرالی ناسازگار دانست
\end{itemize}

\mediumspace

\textbf{نقل‌قول:} \textit{«اگر تمام بشریت یک نظر داشتند و فقط یک نفر نظر مخالف، خاموش کردن او نادرست‌تر از آن نخواهد بود که او، اگر قدرت داشت، بشریت را خاموش کند.»}
\end{personbox}

\begin{defbox}[تعریف: اصل آسیب (Harm Principle)]
اصل بنیادین میل در \book{درباره آزادی}:

\textit{«تنها هدفی که به خاطر آن می‌توان به درستی قدرت را بر عضوی از جامعه متمدن، برخلاف میلش، اعمال کرد، جلوگیری از آسیب به دیگران است.»}

\textbf{پیامدها:}
\begin{itemize}[nosep]
    \item مصرف مواد مخدر؟ (آسیب به خود ≠ آسیب به دیگران)
    \item سخنان نفرت‌انگیز؟ (آیا «آسیب» شامل اهانت می‌شود؟)
    \item فحشا؟ (اگر داوطلبانه باشد؟)
\end{itemize}
این اصل هنوز در بحث‌های امروز حضور دارد.
\end{defbox}

% ══════════════════════════════════════════════════════════════════════════════
\section{لیبرالیسم اجتماعی: چرخش به چپ}
% ══════════════════════════════════════════════════════════════════════════════

در اواخر قرن نوزدهم، برخی لیبرال‌ها دیدند که «آزادی منفی» کافی نیست. فقر و نابرابری، آزادی واقعی را ناممکن می‌کند.

\begin{defbox}[تعریف: لیبرالیسم اجتماعی (نیولیبرالیسم قدیم)]
گرایشی که از اواخر قرن ۱۹ شکل گرفت:
\begin{itemize}[nosep]
    \item \textbf{آزادی مثبت:} آزادی فقط نبود مانع نیست، بلکه توانایی عمل است
    \item کسی که فقیر یا بیمار است، آزاد نیست
    \item دولت می‌تواند و باید برای افزایش آزادی واقعی مداخله کند
    \item دولت رفاه با لیبرالیسم سازگار است
\end{itemize}

\textbf{چهره‌ها:} تی.اچ. گرین، ال.تی. هابهاوس، جان دیویی

\textbf{توجه:} این «نیولیبرالیسم» با «نئولیبرالیسم» دهه ۱۹۸۰ کاملاً متفاوت است!
\end{defbox}

\begin{personbox}[تی.اچ. گرین (۱۸۳۶-۱۸۸۲)]

\textbf{ملیت:} انگلیسی

\textbf{ایده محوری:} آزادی مثبت

\mediumspace

\person{گرین} استدلال کرد:

\begin{itemize}[nosep]
    \item آزادی صرفاً نبود مانع خارجی نیست
    \item آزادی یعنی \textbf{قدرت انجام چیزهای ارزشمند}
    \item فقر، بیماری، جهل—همه موانع آزادی‌اند
    \item دولت می‌تواند با آموزش، بهداشت، قانون کار، آزادی را افزایش دهد
\end{itemize}

\mediumspace

\textbf{تأثیر:} توجیه فکری برای اصلاحات لیبرال‌ها در بریتانیا (۱۹۰۶-۱۹۱۴)
\end{personbox}

\begin{figure}[htbp]
\centering
\begin{tikzpicture}[scale=0.85]

% عنوان
\node[font=\large\bfseries, text=DarkGray] at (0, 6.5) {\fa{تحول لیبرالیسم: از کلاسیک به اجتماعی}};

% محور
\draw[line width=4pt, MediumGray] (-6, 3.5) -- (6, 3.5);

% لیبرالیسم کلاسیک
\fill[RightBlue, opacity=0.7, rounded corners=5pt] (-6, 4) rectangle (-2, 5.5);
\node[text=white, font=\bfseries] at (-4, 5) {\fa{لیبرالیسم کلاسیک}};
\node[text=VeryLightGray, font=\scriptsize] at (-4, 4.4) {\fa{آزادی منفی}};

% لیبرالیسم اجتماعی
\fill[CenterGreen, opacity=0.7, rounded corners=5pt] (2, 4) rectangle (6, 5.5);
\node[text=white, font=\bfseries] at (4, 5) {\fa{لیبرالیسم اجتماعی}};
\node[text=VeryLightGray, font=\scriptsize] at (4, 4.4) {\fa{آزادی مثبت}};

% فلش
\draw[thickarrow, color=GoldAccent, line width=2pt] (-1.8, 4.75) -- (1.8, 4.75);
\node[text=GoldAccent, font=\small] at (0, 5.3) {\fa{اواخر قرن ۱۹}};

% تفاوت‌ها
\node[text=RightBlueDark, font=\scriptsize, text width=4cm, align=center] at (-4, 2.5) {
    \fa{دولت کوچک}\\
    \fa{بازار آزاد}\\
    \fa{فقط حقوق منفی}
};

\node[text=CenterGreenDark, font=\scriptsize, text width=4cm, align=center] at (4, 2.5) {
    \fa{دولت فعال}\\
    \fa{رفاه و آموزش}\\
    \fa{حقوق مثبت هم}
};

\end{tikzpicture}
\caption{تحول از لیبرالیسم کلاسیک به اجتماعی}
\label{fig:liberalism-evolution}
\end{figure}

% ══════════════════════════════════════════════════════════════════════════════
\section{بحران و احیا: نئولیبرالیسم}
% ══════════════════════════════════════════════════════════════════════════════

\subsection{جامعه مون پلرن: نقطه آغاز}

\begin{historybox}[جامعه مون پلرن (۱۹۴۷)]
در ۱۹۴۷، گروهی از اقتصاددانان و فیلسوفان در کوه مون پلرن سوئیس گرد آمدند:

\textbf{بنیان‌گذاران:} هایک، فریدمن، میزس، پوپر، و دیگران

\textbf{هدف:} احیای لیبرالیسم کلاسیک در برابر:
\begin{itemize}[nosep]
    \item کمونیسم شوروی
    \item دولت رفاه کینزی
    \item فاشیسم (تازه شکست‌خورده)
\end{itemize}

\textbf{ایده:} دولت بزرگ—چه سوسیالیستی، چه فاشیستی، چه رفاهی—تهدیدی برای آزادی است.

\textbf{تأثیر:} شبکه‌ای از اندیشکده‌ها که چهار دهه بعد، با تاچر و ریگان، به قدرت رسیدند.
\end{historybox}

\subsection{هایک و راه بردگی}

\begin{personbox}[فردریش آگوست فون هایک (۱۸۹۹-۱۹۹۲)]

\textbf{ملیت:} اتریشی-بریتانیایی

\textbf{آثار کلیدی:} \book{راه بردگی} (۱۹۴۴)، \book{قانون، قانون‌گذاری و آزادی}

\textbf{جایزه نوبل اقتصاد:} ۱۹۷۴

\mediumspace

\person{هایک} مهم‌ترین نظریه‌پرداز نئولیبرالیسم است:

\textbf{استدلال‌های کلیدی:}
\begin{enumerate}[nosep]
    \item \textbf{دانش پراکنده:} هیچ برنامه‌ریزی نمی‌تواند دانش میلیون‌ها فرد را جمع کند
    \item \textbf{قیمت‌ها:} قیمت‌های بازار، این دانش را تجمیع و منتقل می‌کنند
    \item \textbf{راه بردگی:} برنامه‌ریزی مرکزی ضرورتاً به استبداد می‌انجامد
    \item \textbf{نظم خودجوش:} بازار، زبان، قانون—همه محصول «عمل انسان‌اند، نه طراحی»
    \item \textbf{حکومت قانون:} قواعد عام و انتزاعی، نه دستورات خاص
\end{enumerate}

\mediumspace

\textbf{نقل‌قول:} \textit{«کنجکاو است که انسان، با فرض اینکه می‌تواند جامعه را طراحی کند، گمان می‌کند می‌تواند دانشی کسب کند که این کار را ممکن سازد.»}
\end{personbox}

\begin{quotebox}[هایک، راه بردگی]
\textit{«ما به تدریج از مسیری که پیمودیم آگاه شده‌ایم... ما اصول بنیادین تمدن را رها کرده‌ایم... ما در حال انتخاب آگاهانه بین دو نظام اجتماعی نیستیم، بلکه از نظم آزادی به سوی بندگی می‌لغزیم.»}
\end{quotebox}

\begin{debatebox}[آیا هایک محافظه‌کار بود؟]
\textbf{هایک خودش گفت «نه»:}
\begin{itemize}[nosep]
    \item مقاله معروف: «چرا من محافظه‌کار نیستم» (۱۹۶۰)
    \item محافظه‌کاری ضد تغییر است؛ من ضد نیستم
    \item سنت خودبه‌خود ارزش ندارد
    \item من «ویگ قدیمی» یا لیبرال کلاسیک هستم
\end{itemize}

\textbf{اما:}
\begin{itemize}[nosep]
    \item تأکید بر «نظم خودجوش» به برک نزدیک است
    \item نقد «عقل‌گرایی سازنده» محافظه‌کارانه است
    \item در عمل، با محافظه‌کاران متحد شد
\end{itemize}
\end{debatebox}

\subsection{فریدمن و مکتب شیکاگو}

\begin{personbox}[میلتون فریدمن (۱۹۱۲-۲۰۰۶)]

\textbf{ملیت:} آمریکایی

\textbf{آثار کلیدی:} \book{سرمایه‌داری و آزادی} (۱۹۶۲)، \book{آزاد برای انتخاب} (۱۹۸۰)

\textbf{جایزه نوبل اقتصاد:} ۱۹۷۶

\mediumspace

\person{فریدمن} نئولیبرالیسم را عامه‌فهم کرد:

\textbf{ایده‌های کلیدی:}
\begin{itemize}[nosep]
    \item \textbf{مانیتاریسم:} تورم همیشه پدیده‌ای پولی است
    \item \textbf{نقد کینز:} سیاست مالی بی‌اثر یا مضر است
    \item \textbf{کوپن آموزشی:} انتخاب مدرسه توسط والدین
    \item \textbf{مالیات منفی:} جایگزین رفاه بوروکراتیک
    \item \textbf{ارتش داوطلب:} پایان سربازی اجباری
    \item \textbf{آزادسازی مواد مخدر:} جنگ با مواد شکست‌خورده
\end{itemize}

\mediumspace

\textbf{نقل‌قول:} \textit{«هیچ‌چیز به اندازه برنامه دولتی ماندگار نیست که قرار بود موقت باشد.»}
\end{personbox}

% ══════════════════════════════════════════════════════════════════════════════
\section{نئولیبرالیسم در عمل}
% ══════════════════════════════════════════════════════════════════════════════

\begin{figure}[htbp]
\centering
\begin{tikzpicture}[scale=0.85]

% عنوان
\node[font=\large\bfseries, text=DarkGray] at (0, 7.5) {\fa{تایم‌لاین: عروج نئولیبرالیسم}};

% خط زمان
\draw[line width=4pt, MediumGray] (-7, 5.5) -- (7, 5.5);

% رویدادها
% مون پلرن
\fill[MediumGray] (-6, 5.5) circle (8pt);
\node[above, text=MediumGray, font=\scriptsize\bfseries] at (-6, 6) {\fa{۱۹۴۷}};
\node[below, text width=2cm, align=center, text=DarkGray, font=\tiny] at (-6, 4.9) {
    \fa{جامعه}\\
    \fa{مون پلرن}
};

% شیلی
\fill[OrangeAccent] (-3.5, 5.5) circle (9pt);
\node[above, text=OrangeAccent, font=\scriptsize\bfseries] at (-3.5, 6) {\fa{۱۹۷۳}};
\node[below, text width=2cm, align=center, text=DarkGray, font=\tiny] at (-3.5, 4.9) {
    \fa{شیلی پینوشه}\\
    \fa{«پسران شیکاگو»}
};

% تاچر
\fill[RightBlue] (-1, 5.5) circle (10pt);
\node[above, text=RightBlue, font=\scriptsize\bfseries] at (-1, 6.1) {\fa{۱۹۷۹}};
\node[below, text width=2cm, align=center, text=DarkGray, font=\tiny] at (-1, 4.8) {
    \fa{تاچر}\\
    \fa{بریتانیا}
};

% ریگان
\fill[RightBlue] (1.5, 5.5) circle (10pt);
\node[above, text=RightBlue, font=\scriptsize\bfseries] at (1.5, 6.1) {\fa{۱۹۸۰}};
\node[below, text width=2cm, align=center, text=DarkGray, font=\tiny] at (1.5, 4.8) {
    \fa{ریگان}\\
    \fa{آمریکا}
};

% اجماع واشنگتن
\fill[GoldAccent] (4, 5.5) circle (9pt);
\node[above, text=GoldAccent, font=\scriptsize\bfseries] at (4, 6) {\fa{۱۹۸۹}};
\node[below, text width=2.2cm, align=center, text=DarkGray, font=\tiny] at (4, 4.8) {
    \fa{اجماع واشنگتن}\\
    \fa{جهانی‌شدن}
};

% بحران
\fill[LeftRed] (6.5, 5.5) circle (9pt);
\node[above, text=LeftRed, font=\scriptsize\bfseries] at (6.5, 6) {\fa{۲۰۰۸}};
\node[below, text width=2cm, align=center, text=DarkGray, font=\tiny] at (6.5, 4.9) {
    \fa{بحران مالی}\\
    \fa{زیر سؤال}
};

\end{tikzpicture}
\caption{سیر عروج نئولیبرالیسم}
\label{fig:neoliberalism-timeline}
\end{figure}

\subsection{شیلی: آزمایشگاه نئولیبرال}

\begin{historybox}[شیلی پینوشه: اولین آزمایش]
در ۱۹۷۳، ژنرال پینوشه با کودتا علیه آلنده سوسیالیست به قدرت رسید:
\begin{itemize}[nosep]
    \item «پسران شیکاگو»: اقتصاددانان شیلیایی تحصیل‌کرده نزد فریدمن
    \item سیاست‌ها: خصوصی‌سازی، آزادسازی، کاهش رفاه
    \item نتایج اولیه: تورم بالا، بیکاری، فقر
    \item نتایج بلندمدت: رشد اقتصادی، اما نابرابری شدید
\end{itemize}

\textbf{مناقشه:} آیا نئولیبرالیسم به دیکتاتوری نیاز داشت؟

\textbf{هایک (۱۹۸۱):} «دیکتاتوری موقت ممکن است لازم باشد تا آزادی بلندمدت ممکن شود.»
\end{historybox}

\subsection{اجماع واشنگتن}

\begin{defbox}[تعریف: اجماع واشنگتن (۱۹۸۹)]
اصطلاحی که جان ویلیامسون برای سیاست‌های توصیه‌شده به کشورهای در حال توسعه به کار برد:

\textbf{ده دستور:}
\begin{enumerate}[nosep]
    \item انضباط مالی (کسری بودجه پایین)
    \item اولویت‌بندی هزینه‌ها (آموزش، بهداشت، زیرساخت)
    \item اصلاح مالیاتی (نرخ پایین، پایه گسترده)
    \item آزادسازی نرخ بهره
    \item نرخ ارز رقابتی
    \item آزادسازی تجاری
    \item آزادسازی سرمایه‌گذاری خارجی
    \item خصوصی‌سازی
    \item مقررات‌زدایی
    \item حقوق مالکیت
\end{enumerate}

\textbf{اجراکنندگان:} صندوق بین‌المللی پول، بانک جهانی

\textbf{نقد:} این سیاست‌ها در بسیاری کشورها (آرژانتین، روسیه) به بحران انجامید.
\end{defbox}

% ══════════════════════════════════════════════════════════════════════════════
\section{نقد نئولیبرالیسم}
% ══════════════════════════════════════════════════════════════════════════════

\subsection{نقد از چپ}

\begin{itemize}[nosep, rightmargin=1.5cm]
    \item نئولیبرالیسم = \textbf{سرمایه‌داری افسارگسیخته}
    \item افزایش نابرابری در همه کشورها
    \item تضعیف اتحادیه‌ها و حقوق کارگران
    \item بحران مالی ۲۰۰۸ نتیجه مقررات‌زدایی بود
    \item کالایی‌سازی همه چیز: آموزش، بهداشت، آب، هوا
    \item «آزادی» فقط برای ثروتمندان
\end{itemize}

\subsection{نقد از راست محافظه‌کار}

\begin{itemize}[nosep, rightmargin=1.5cm]
    \item نئولیبرالیسم جوامع و سنت‌ها را ویران کرده
    \item «جامعه وجود ندارد» غلط است—جامعه مهم است
    \item جهانی‌شدن به ضرر کارگران بومی است
    \item فقط «آزادی اقتصادی» کافی نیست—ارزش‌ها هم مهم‌اند
    \item خانواده و ملت فدای بازار شده‌اند
\end{itemize}

\begin{debatebox}[بحران ۲۰۰۸: پایان نئولیبرالیسم؟]
\textbf{بله:}
\begin{itemize}[nosep]
    \item بحران نشان داد بازار خودتنظیم نیست
    \item دولت‌ها مجبور شدند بانک‌ها را نجات دهند
    \item «پایان تاریخ» تمام شد
\end{itemize}

\textbf{نه:}
\begin{itemize}[nosep]
    \item پس از بحران، همان سیاست‌ها ادامه یافت (ریاضت اقتصادی)
    \item جایگزینی ارائه نشد
    \item نئولیبرالیسم «زامبی» شده—مرده اما راه می‌رود
\end{itemize}
\end{debatebox}

% ══════════════════════════════════════════════════════════════════════════════
\section{لیبرالیسم معاصر: رالز و منتقدانش}
% ══════════════════════════════════════════════════════════════════════════════

\begin{personbox}[جان رالز (۱۹۲۱-۲۰۰۲)]

\textbf{ملیت:} آمریکایی

\textbf{اثر کلیدی:} \book{نظریه‌ای در باب عدالت} (۱۹۷۱)

\textbf{ایده محوری:} عدالت به مثابه انصاف

\mediumspace

\person{رالز} فلسفه سیاسی را احیا کرد و لیبرالیسمی جدید ارائه داد:

\textbf{موقعیت اصیل و پرده جهل:}
\begin{itemize}[nosep]
    \item فرض کنید نمی‌دانید در جامعه چه جایگاهی خواهید داشت
    \item پشت این «پرده جهل»، چه اصولی انتخاب می‌کنید؟
\end{itemize}

\textbf{دو اصل عدالت:}
\begin{enumerate}[nosep]
    \item \textbf{اصل آزادی:} بیشترین آزادی‌های پایه برابر برای همه
    \item \textbf{اصل تفاوت:} نابرابری فقط وقتی موجه است که به نفع محروم‌ترین‌ها باشد
\end{enumerate}

\mediumspace

\textbf{پیامد:} دفاع فلسفی از دولت رفاه لیبرال
\end{personbox}

\begin{comparebox}[مقایسه: لیبرالیسم رالز و نئولیبرالیسم]
\begin{table}[H]
\centering
\begin{tabular}{r|c|c}
\toprule
\textbf{موضوع} & \textbf{رالز} & \textbf{نئولیبرالیسم} \\
\midrule
عدالت & توزیعی (نتیجه مهم است) & رویه‌ای (فرآیند مهم است) \\
نابرابری & باید توجیه شود & نتیجه طبیعی بازار \\
دولت & رفاهی، فعال & حداقلی \\
بازار & ابزار، نه هدف & تقریباً مقدس \\
آزادی & منفی و مثبت & عمدتاً منفی \\
\bottomrule
\end{tabular}
\end{table}
\end{comparebox}

% ══════════════════════════════════════════════════════════════════════════════
\section{اختلافات درون لیبرالیسم}
% ══════════════════════════════════════════════════════════════════════════════

\begin{figure}[htbp]
\centering
\begin{tikzpicture}[scale=0.85]

% عنوان
\node[font=\large\bfseries, text=DarkGray] at (0, 6.5) {\fa{طیف درونی لیبرالیسم}};

% محور
\draw[line width=4pt, MediumGray] (-6, 3.5) -- (6, 3.5);
\node[text=LeftRed, font=\small\bfseries] at (-6.5, 3.5) {\fa{چپ}};
\node[text=RightBlue, font=\small\bfseries] at (6.5, 3.5) {\fa{راست}};

% نقاط
% لیبرالیسم برابری‌خواه
\fill[LeftRedLight] (-4, 3.5) circle (10pt);
\node[above, text=LeftRedDark, font=\small\bfseries] at (-4, 4) {\fa{رالز}};
\node[below, text=DarkGray, font=\scriptsize, text width=2.5cm, align=center] at (-4, 2.8) {
    \fa{برابری‌خواه}\\
    \fa{دولت رفاه}
};

% لیبرالیسم کلاسیک
\fill[CenterGreen] (0, 3.5) circle (10pt);
\node[above, text=CenterGreenDark, font=\small\bfseries] at (0, 4) {\fa{میل}};
\node[below, text=DarkGray, font=\scriptsize, text width=2.5cm, align=center] at (0, 2.8) {
    \fa{کلاسیک}\\
    \fa{دولت محدود}
};

% نئولیبرالیسم
\fill[RightBlueLight] (3, 3.5) circle (10pt);
\node[above, text=RightBlueDark, font=\small\bfseries] at (3, 4) {\fa{هایک}};
\node[below, text=DarkGray, font=\scriptsize, text width=2.5cm, align=center] at (3, 2.8) {
    \fa{نئولیبرال}\\
    \fa{بازار آزاد}
};

% لیبرتارینیسم
\fill[RightBlueDark] (5.5, 3.5) circle (10pt);
\node[above, text=RightBlueDark, font=\small\bfseries] at (5.5, 4) {\fa{نوزیک}};
\node[below, text=DarkGray, font=\scriptsize, text width=2.5cm, align=center] at (5.5, 2.8) {
    \fa{لیبرتارین}\\
    \fa{دولت حداقلی}
};

\end{tikzpicture}
\caption{طیف گرایش‌های درون لیبرالیسم}
\label{fig:liberalism-spectrum}
\end{figure}

\begin{debatebox}[اختلاف کلیدی: دولت چقدر فعال باشد؟]
\textbf{لیبرال‌های چپ (رالز، دورکین):}
\begin{itemize}[nosep]
    \item آزادی واقعی نیاز به منابع دارد
    \item دولت باید فرصت‌ها را برابر کند
    \item بهداشت و آموزش حق است
    \item بازتوزیع، آزادی را افزایش می‌دهد
\end{itemize}

\textbf{لیبرال‌های راست (هایک، نوزیک):}
\begin{itemize}[nosep]
    \item دولت بزرگ، تهدید آزادی است
    \item بازتوزیع = تصاحب اجباری = نوعی بردگی
    \item بازار بهترین توزیع‌کننده است
    \item «حقوق مثبت» توهم است
\end{itemize}
\end{debatebox}

% ══════════════════════════════════════════════════════════════════════════════
%                         خلاصه و پرسش‌ها
% ══════════════════════════════════════════════════════════════════════════════

\newpage

\begin{summarybox}

\textbf{۱. سردرگمی:} «لیبرال» در آمریکا (چپ) و اروپا (راست اقتصادی) معنای متفاوتی دارد.

\textbf{۲. هسته مشترک:} فردگرایی، آزادی، برابری، عقل‌گرایی، حکومت محدود.

\textbf{۳. کلاسیک:} لاک (حقوق طبیعی)، اسمیت (بازار)، میل (آزادی و اصل آسیب).

\textbf{۴. اجتماعی:} گرین و دیگران آزادی مثبت را افزودند—توجیه دولت رفاه.

\textbf{۵. نئولیبرالیسم:} هایک و فریدمن لیبرالیسم کلاسیک را احیا کردند. مون پلرن، تاچر، ریگان.

\textbf{۶. در عمل:} شیلی، اجماع واشنگتن، جهانی‌شدن.

\textbf{۷. نقد:} از چپ (نابرابری) و از راست محافظه‌کار (تخریب جامعه).

\textbf{۸. رالز:} عدالت به مثابه انصاف. لیبرالیسم برابری‌خواه.

\textbf{۹. اختلافات:} کلاسیک vs اجتماعی vs نئولیبرال vs لیبرتارین.

\end{summarybox}

\vspace{20pt}

\begin{questionbox}

\begin{enumerate}[nosep, rightmargin=1cm]
    \item آیا اسمیت واقعاً «بنیادگرای بازار» بود؟ یا نئولیبرال‌ها او را سوءتفسیر کرده‌اند؟

    \item اصل آسیب میل می‌گوید فقط آسیب به دیگران توجیه محدودیت است. آیا «آسیب» شامل اهانت، تنفر یا آسیب روانی هم می‌شود؟

    \item آیا نئولیبرالیسم به دیکتاتوری نیاز داشت؟ (شیلی، چین؟)

    \item رالز و هایک هر دو «لیبرال»اند، اما نتیجه‌گیری‌هایشان کاملاً متفاوت است. کدام‌یک به سنت لیبرالی وفادارتر است؟
\end{enumerate}

\end{questionbox}

% ══════════════════════════════════════════════════════════════════════════════
%                         پایان فصل دهم
% ══════════════════════════════════════════════════════════════════════════════